% Options for packages loaded elsewhere
\PassOptionsToPackage{unicode}{hyperref}
\PassOptionsToPackage{hyphens}{url}
\PassOptionsToPackage{dvipsnames,svgnames,x11names}{xcolor}
%
\documentclass[
  10pt,
]{article}
\usepackage{amsmath,amssymb}
\usepackage{iftex}
\ifPDFTeX
  \usepackage[T1]{fontenc}
  \usepackage[utf8]{inputenc}
  \usepackage{textcomp} % provide euro and other symbols
\else % if luatex or xetex
  \usepackage{unicode-math} % this also loads fontspec
  \defaultfontfeatures{Scale=MatchLowercase}
  \defaultfontfeatures[\rmfamily]{Ligatures=TeX,Scale=1}
\fi
\usepackage{lmodern}
\ifPDFTeX\else
  % xetex/luatex font selection
  \setmainfont[]{DejaVu Serif}
  \setsansfont[]{DejaVu Sans}
  \setmonofont[]{DejaVu Sans Mono}
  \setmathfont[]{Latin Modern Math}
\fi
% Use upquote if available, for straight quotes in verbatim environments
\IfFileExists{upquote.sty}{\usepackage{upquote}}{}
\IfFileExists{microtype.sty}{% use microtype if available
  \usepackage[]{microtype}
  \UseMicrotypeSet[protrusion]{basicmath} % disable protrusion for tt fonts
}{}
\makeatletter
\@ifundefined{KOMAClassName}{% if non-KOMA class
  \IfFileExists{parskip.sty}{%
    \usepackage{parskip}
  }{% else
    \setlength{\parindent}{0pt}
    \setlength{\parskip}{6pt plus 2pt minus 1pt}}
}{% if KOMA class
  \KOMAoptions{parskip=half}}
\makeatother
\usepackage{xcolor}
\usepackage[margin=2.0cm]{geometry}
\usepackage{color}
\usepackage{fancyvrb}
\newcommand{\VerbBar}{|}
\newcommand{\VERB}{\Verb[commandchars=\\\{\}]}
\DefineVerbatimEnvironment{Highlighting}{Verbatim}{commandchars=\\\{\}}
% Add ',fontsize=\small' for more characters per line
\newenvironment{Shaded}{}{}
\newcommand{\AlertTok}[1]{\textcolor[rgb]{1.00,0.00,0.00}{\textbf{#1}}}
\newcommand{\AnnotationTok}[1]{\textcolor[rgb]{0.38,0.63,0.69}{\textbf{\textit{#1}}}}
\newcommand{\AttributeTok}[1]{\textcolor[rgb]{0.49,0.56,0.16}{#1}}
\newcommand{\BaseNTok}[1]{\textcolor[rgb]{0.25,0.63,0.44}{#1}}
\newcommand{\BuiltInTok}[1]{\textcolor[rgb]{0.00,0.50,0.00}{#1}}
\newcommand{\CharTok}[1]{\textcolor[rgb]{0.25,0.44,0.63}{#1}}
\newcommand{\CommentTok}[1]{\textcolor[rgb]{0.38,0.63,0.69}{\textit{#1}}}
\newcommand{\CommentVarTok}[1]{\textcolor[rgb]{0.38,0.63,0.69}{\textbf{\textit{#1}}}}
\newcommand{\ConstantTok}[1]{\textcolor[rgb]{0.53,0.00,0.00}{#1}}
\newcommand{\ControlFlowTok}[1]{\textcolor[rgb]{0.00,0.44,0.13}{\textbf{#1}}}
\newcommand{\DataTypeTok}[1]{\textcolor[rgb]{0.56,0.13,0.00}{#1}}
\newcommand{\DecValTok}[1]{\textcolor[rgb]{0.25,0.63,0.44}{#1}}
\newcommand{\DocumentationTok}[1]{\textcolor[rgb]{0.73,0.13,0.13}{\textit{#1}}}
\newcommand{\ErrorTok}[1]{\textcolor[rgb]{1.00,0.00,0.00}{\textbf{#1}}}
\newcommand{\ExtensionTok}[1]{#1}
\newcommand{\FloatTok}[1]{\textcolor[rgb]{0.25,0.63,0.44}{#1}}
\newcommand{\FunctionTok}[1]{\textcolor[rgb]{0.02,0.16,0.49}{#1}}
\newcommand{\ImportTok}[1]{\textcolor[rgb]{0.00,0.50,0.00}{\textbf{#1}}}
\newcommand{\InformationTok}[1]{\textcolor[rgb]{0.38,0.63,0.69}{\textbf{\textit{#1}}}}
\newcommand{\KeywordTok}[1]{\textcolor[rgb]{0.00,0.44,0.13}{\textbf{#1}}}
\newcommand{\NormalTok}[1]{#1}
\newcommand{\OperatorTok}[1]{\textcolor[rgb]{0.40,0.40,0.40}{#1}}
\newcommand{\OtherTok}[1]{\textcolor[rgb]{0.00,0.44,0.13}{#1}}
\newcommand{\PreprocessorTok}[1]{\textcolor[rgb]{0.74,0.48,0.00}{#1}}
\newcommand{\RegionMarkerTok}[1]{#1}
\newcommand{\SpecialCharTok}[1]{\textcolor[rgb]{0.25,0.44,0.63}{#1}}
\newcommand{\SpecialStringTok}[1]{\textcolor[rgb]{0.73,0.40,0.53}{#1}}
\newcommand{\StringTok}[1]{\textcolor[rgb]{0.25,0.44,0.63}{#1}}
\newcommand{\VariableTok}[1]{\textcolor[rgb]{0.10,0.09,0.49}{#1}}
\newcommand{\VerbatimStringTok}[1]{\textcolor[rgb]{0.25,0.44,0.63}{#1}}
\newcommand{\WarningTok}[1]{\textcolor[rgb]{0.38,0.63,0.69}{\textbf{\textit{#1}}}}
\usepackage{longtable,booktabs,array}
\usepackage{calc} % for calculating minipage widths
% Correct order of tables after \paragraph or \subparagraph
\usepackage{etoolbox}
\makeatletter
\patchcmd\longtable{\par}{\if@noskipsec\mbox{}\fi\par}{}{}
\makeatother
% Allow footnotes in longtable head/foot
\IfFileExists{footnotehyper.sty}{\usepackage{footnotehyper}}{\usepackage{footnote}}
\makesavenoteenv{longtable}
\setlength{\emergencystretch}{3em} % prevent overfull lines
\providecommand{\tightlist}{%
  \setlength{\itemsep}{0pt}\setlength{\parskip}{0pt}}
\setcounter{secnumdepth}{-\maxdimen} % remove section numbering
% doc_header.tex -- LaTeX preamble additions used by build_docs.sh for every
% Markdown -> LaTeX -> PDF conversion (the documents of the folder Revision_Lord_Gravity).
\usepackage{etoolbox}
% longtable may not be loaded by pandoc when a document has no table
\usepackage{longtable}
% bold vectors: \bm -> \symbf under unicode-math (XeLaTeX), bm package otherwise
\makeatletter
\@ifpackageloaded{unicode-math}{\providecommand{\bm}[1]{\symbf{#1}}}{\usepackage{bm}}
\makeatother
% smaller type in long tables so that wide audit tables fit the text width;
% inside a table a line break is allowed after an underscore of plain text
% (pandoc writes it as \_), so that long identifiers wrap inside their column
% rows are set a little apart, so that symbols with indices do not touch
\AtBeginEnvironment{longtable}{\footnotesize\renewcommand{\arraystretch}{1.45}%
  \let\sptunderscore\_\def\_{\sptunderscore\allowbreak}}
\setlength{\LTcapwidth}{\textwidth}
% A short table is kept on one page.  The build (code/scripts/keep_tables.py) ends the rows of such a table with \\*,
% which forbids a page break after a row, and sets \sptkeptpenalty to 10000 in front of the table.  This penalty is
% put here in the place of the breakpoint (penalty 0) that longtable leaves after every table, so that the paragraph
% that follows a kept table, its caption, begins on the same page; it is set back to 0 after the table.
\newcount\sptkeptpenalty
\sptkeptpenalty=0
\makeatletter
\patchcmd{\endlongtable}{\endgraf\penalty\z@\addvspace\LTpost}%
  {\endgraf\penalty\sptkeptpenalty\addvspace\LTpost\global\sptkeptpenalty=\z@}%
  {}{\PackageError{doc_header}{the end of longtable could not be patched}{The tables can still be typeset, but a short table may be separated from its caption.}}
\makeatother
% allow line breaks after slashes and hyphens inside table cells
\usepackage{xurl}
% code blocks: break over-long lines instead of running into the margin.
% Plain ``` blocks are typeset by pandoc as {verbatim}, highlighted ones
% (```bash) as {Highlighting}; both are re-declared as breakable fancyvrb
% environments.  The continuation line is marked with a small hooked arrow.
\usepackage{fvextra}
\RecustomVerbatimEnvironment{verbatim}{Verbatim}{breaklines,breakanywhere,fontsize=\small}
\makeatletter
\@ifundefined{Highlighting}{}{%
  \RecustomVerbatimEnvironment{Highlighting}{Verbatim}{breaklines,breakanywhere,fontsize=\small,commandchars=\\\{\}}}
\makeatother
% inline code (pandoc writes it as \texttt{...}): never hyphenate -- an inserted
% hyphen could be read as part of a file name -- but allow a line break after
% an underscore, slash, comma, colon or equals sign, so that long paths and
% column lists wrap inside the text width and inside table columns.
% Inside a formula (e.g. the name of a notebook symbol in an equation) the same
% command must first leave math mode, which \text of amsmath does.
\ExplSyntaxOn
\tl_new:N \l__spt_code_tl
\cs_new_protected:Npn \__spt_texttt:n #1
  {
    \tl_set:Nn \l__spt_code_tl {#1}
    \regex_replace_all:nnN { ([\/\,\:\=]) } { \1 \c{allowbreak} } \l__spt_code_tl
    \regex_replace_all:nnN { \c{_} } { \c{_} \c{allowbreak} } \l__spt_code_tl
    \group_begin: \ttfamily \hyphenchar \font = -1 \scan_stop: \l__spt_code_tl \group_end:
  }
\RenewDocumentCommand \texttt { m }
  {
    \mode_if_math:TF { \text { \__spt_texttt:n {#1} } } { \__spt_texttt:n {#1} }
  }
\ExplSyntaxOff
% A figure is written in the Markdown source as an image paragraph followed by
% a caption paragraph ("**Figure S1.** ..."), so that the caption is also
% visible on GitHub.  An image that starts a paragraph is centered.
\usepackage{graphicx}% pandoc loads it only for documents that contain an image
\NewCommandCopy\sptincludegraphics\includegraphics
\RenewDocumentCommand\includegraphics{O{} m}{%
  \ifvmode\begin{center}\sptincludegraphics[#1]{#2}\end{center}%
  \else\sptincludegraphics[#1]{#2}\fi}
% keep figures where they are referenced
\usepackage{float}
\floatplacement{figure}{H}
% running header with the document title
\usepackage{fancyhdr}
\pagestyle{fancy}
\fancyhf{}
\fancyhead[L]{\small PROJECT spacetimes, Revision\_Lord\_Gravity}
\fancyhead[R]{\small\thepage}
\renewcommand{\headrulewidth}{0.4pt}
% 16x16 matrices: tighter columns so that a full matrix fits the text width
\setlength{\arraycolsep}{3.2pt}
% The documents of this project contain long tagged symbols, file paths and
% identifiers that cannot be hyphenated.  A paragraph is therefore allowed to
% have looser lines before TeX gives up and lets a line stick out of the margin.
\AtBeginDocument{\tolerance=2000 \emergencystretch=5em \hbadness=2000}
\ifLuaTeX
  \usepackage{selnolig}  % disable illegal ligatures
\fi
\IfFileExists{bookmark.sty}{\usepackage{bookmark}}{\usepackage{hyperref}}
\IfFileExists{xurl.sty}{\usepackage{xurl}}{} % add URL line breaks if available
\urlstyle{same}
\hypersetup{
  pdftitle={The brochon and E. A. Lord's gauge theory of gravitation},
  pdfauthor={Prepared for Patrick Lee Nash, Ph.D.},
  colorlinks=true,
  linkcolor={blue},
  filecolor={Maroon},
  citecolor={blue},
  urlcolor={blue},
  pdfcreator={LaTeX via pandoc}}

\title{The brochon and E. A. Lord's gauge theory of gravitation}
\usepackage{etoolbox}
\makeatletter
\providecommand{\subtitle}[1]{% add subtitle to \maketitle
  \apptocmd{\@title}{\par {\large #1 \par}}{}{}
}
\makeatother
\subtitle{PROJECT spacetimes, folder Revision\_Lord\_Gravity: the
manuscript of task C23}
\author{Prepared for Patrick Lee Nash, Ph.D.}
\date{November 22, 2026}

\begin{document}
\maketitle

\hypertarget{what-this-manuscript-contains-and-the-state-of-the-work}{%
\section{0. What this manuscript contains, and the state of the
work}\label{what-this-manuscript-contains-and-the-state-of-the-work}}

The task of this manuscript is to couple the spinor field of the
project, the \textbf{brochon}, to gravitation by the gauge theory of E.
A. Lord, to solve its field equations in the Schwarzschild and the
Vaidya space-times on both sides of the horizon, and to derive
predictions involving the brochon and the zeugon that gravitation can
detect. It is written in parts, each pushed when it is complete:

\begin{itemize}
\tightlist
\item
  \textbf{Part I} (Sections 2 to 8, this version): Lord's paper,
  ``General relativity from gauge invariance'', Proceedings of the
  Cambridge Philosophical Society 69, 423 (1971), read and restated
  section by section in my words, with every equation that the later
  parts use re-derived and checked by programs. Section 7 lists the
  equations whose printed form does not follow from Lord's definitions,
  with the corrected form and the check that decides.
\item
  \textbf{Part II} (planned): the association of the group
  \(SL(2,\mathbb C)\) of the project with Lord's gauging of the
  representation of the Clifford algebra.
\item
  \textbf{Parts III to VI} (planned): the gravitational coupling of the
  brochon by Lord's covariant derivative; the brochon in the
  Schwarzschild field; the brochon in the Vaidya metric; the brochon
  pair bound by the zeugon and the predictions.
\end{itemize}

Section 9 gives the provenance of Part I: the three programs that check
it and their complete reports. The checks of Part II are already made by
the same programs and are printed there too.

\hypertarget{conventions-units-and-notation}{%
\section{1. Conventions, units and
notation}\label{conventions-units-and-notation}}

Units with \(\hbar=c=1\). Space-time indices
\(\mu,\nu,\rho,\sigma,\dots\) run over \(0,1,2,3\); the Minkowski metric
is \(\eta^{\rm LL}=\mathrm{diag}(1,-1,-1,-1)\), written with indices as
\(\eta^{\rm LL}_{\mu\nu}\) and its inverse \(\eta^{{\rm LL}\,\mu\nu}\)
(the same numbers). A repeated index is summed. The suffixes mark the
origin of a symbol: LL for the objects of the earlier work on the
Lorentz super-loop, SF for those of the earlier hypotheses of the
project, LG for those that this folder introduces, and \textbf{Lord} for
the objects of Lord's paper that the project does not otherwise have.
Where an object of Lord's paper is the same as an object of the project,
the project's symbol is used; Part II proves each such identity (checks
P-1 to P-4 of Section 9).

\textbf{The matrices of Lord's Section 1.} The four \(2\times2\)
matrices \(\sigma^{\rm LL}_\mu\) are the unit matrix and the Pauli
matrices, \(\sigma^{\rm LL}_0=\mathsf 1_2\),
\(\sigma^{\rm LL}_1=\begin{pmatrix}0&1\\1&0\end{pmatrix}\),
\(\sigma^{\rm LL}_2=\begin{pmatrix}0&-i\\i&0\end{pmatrix}\),
\(\sigma^{\rm LL}_3=\begin{pmatrix}1&0\\0&-1\end{pmatrix}\); Lord writes
them with an upper index, \(\sigma^\mu=(1,\boldsymbol\sigma)\). The
matrix \(\mathsf E_{\rm LL}=\begin{pmatrix}0&1\\-1&0\end{pmatrix}\) is
his \(\epsilon\). For any \(2\times2\) matrix \(X\) the \textbf{bar} is
\[\bar X=-\mathsf E_{\rm LL}X^T\mathsf E_{\rm LL},\tag{1.1}\] which is
the adjugate of \(X\) (the transpose of its matrix of cofactors), so
that \(X\bar X=\bar XX=\det X\,\mathsf 1_2\); for the Pauli set it gives
\(\bar\sigma^{\rm LL}_0=\mathsf 1_2\) and
\(\bar\sigma^{\rm LL}_j=-\sigma^{\rm LL}_j\). The \textbf{Dirac matrices
in the Weyl form} (the chiral basis of the project) are, in \(2\times2\)
blocks,
\[\gamma^\mu_{\rm LL}=\begin{pmatrix}0&\sigma^{\rm LL}_\mu\\\bar\sigma^{\rm LL}_\mu&0\end{pmatrix},\qquad
\gamma^{\rm ch}_{\rm LL}=i\gamma^0_{\rm LL}\gamma^1_{\rm LL}\gamma^2_{\rm LL}\gamma^3_{\rm LL}=\begin{pmatrix}-\mathsf 1_2&0\\0&\mathsf 1_2\end{pmatrix},\qquad
C_{\rm LL}=\begin{pmatrix}\mathsf E_{\rm LL}&0\\0&-\mathsf E_{\rm LL}\end{pmatrix}.\tag{1.2}\]
Lord's chirality matrix is
\(\gamma^5=-i\gamma^0\gamma^1\gamma^2\gamma^3=-\gamma^{\rm ch}_{\rm LL}\);
it is written \(-\gamma^{\rm ch}_{\rm LL}\) here. The products are
\(\gamma^{\mu\nu}_{\rm LL}=\tfrac12(\gamma^\mu_{\rm LL}\gamma^\nu_{\rm LL}-\gamma^\nu_{\rm LL}\gamma^\mu_{\rm LL})\),
and
\(\sigma^{\mu\nu}_{\rm LG}=\tfrac12(\sigma^{\rm LL}_\mu\bar\sigma^{\rm LL}_\nu-\sigma^{\rm LL}_\nu\bar\sigma^{\rm LL}_\mu)\)
(Lord's \(\sigma^{\mu\nu}\)). The Levi-Civita symbol
\(\epsilon^{\rm LG}_{\mu\nu\rho\sigma}\) is totally antisymmetric with
\(\epsilon^{\rm LG}_{0123}=+1\); the \textbf{dual} of an antisymmetric
tensor is
\(f_{\bar\mu\bar\nu}=\tfrac12i\,\epsilon^{\rm LG}_{\mu\nu\rho\sigma}f^{\rho\sigma}\)
(Lord's (1.6)).

\textbf{Traces.} \(\mathrm{tr}\) is the trace of a matrix;
\(\mathrm{tr}_{\rm N}=\tfrac14\mathrm{tr}\) is the trace of a
\(4\times4\) matrix normalized so that
\(\mathrm{tr}_{\rm N}\mathsf 1_4=1\). Section 5 shows that Lord's ``tr''
in his Section 4 must be read as \(\mathrm{tr}_{\rm N}\).

\textbf{Equation numbers.} An equation number without a name, such as
(5.3), refers to this manuscript; Lord's equations are cited as ``Lord
(4.13)''. Lord's paper is not in the repository; the copy that was read
is identified by its SHA-256 checksum in Section 9.

\hypertarget{part-i.-lords-paper-read-and-understood}{%
\section{Part I. Lord's paper, read and
understood}\label{part-i.-lords-paper-read-and-understood}}

Lord's paper has five sections. The first sets out the free field
equations of particles of any spin in the form of Bargmann and Wigner,
and their two invariance groups; the second shows that the massless
equations of spin 1 and 2 are Maxwell's equations and the linear spin-2
equations, with potentials; the third lets the spin-2 field interact;
the fourth makes the group of the changes of representation local, which
forces auxiliary fields into the derivative; the fifth identifies those
auxiliary fields with the potentials of the spin-2 field, and shows that
the resulting non-linear theory is Einstein's general relativity.
Sections 2 to 6 below follow his five sections.

\hypertarget{the-free-fields-and-their-two-invariance-groups-lords-section-1}{%
\section{2. The free fields and their two invariance groups (Lord's
Section
1)}\label{the-free-fields-and-their-two-invariance-groups-lords-section-1}}

\hypertarget{two-component-spinors}{%
\subsection{2.1 Two-component spinors}\label{two-component-spinors}}

The matrices of Section 1 obey the relations that make them a basis of
the Pauli algebra (a representation of the quaternions):
\[\tfrac12\big(\sigma^{\rm LL}_\mu\bar\sigma^{\rm LL}_\nu+\sigma^{\rm LL}_\nu\bar\sigma^{\rm LL}_\mu\big)=\eta^{{\rm LL}\,\mu\nu}\mathsf 1_2,\tag{2.1}\]
\[[\sigma^{\mu\nu}_{\rm LG},\sigma^{\rho\sigma}_{\rm LG}]=2\big(\eta^{{\rm LL}\,\nu\rho}\sigma^{\mu\sigma}_{\rm LG}-\eta^{{\rm LL}\,\mu\rho}\sigma^{\nu\sigma}_{\rm LG}+\eta^{{\rm LL}\,\mu\sigma}\sigma^{\nu\rho}_{\rm LG}-\eta^{{\rm LL}\,\nu\sigma}\sigma^{\mu\rho}_{\rm LG}\big)\tag{2.2}\]
(Lord (1.3), (1.5)), so that the six matrices
\(\tfrac12\sigma^{\mu\nu}_{\rm LG}\) generate a representation of the
homogeneous Lorentz group; they are self-dual,
\(\sigma_{\mu\nu\,{\rm LG}}=\sigma_{\bar\mu\bar\nu\,{\rm LG}}\) (Lord
(1.7)). A Lorentz transformation \(\Lambda_{{\rm Lord}\,\nu}{}^\mu\)
(\(\Lambda_{{\rm Lord}\,\mu}{}^\rho\Lambda_{{\rm Lord}\,\lambda}{}^\sigma\eta^{\rm LL}_{\rho\sigma}=\eta^{\rm LL}_{\mu\lambda}\))
is represented on two-component spinors by a matrix \(L_{\rm Lord}\) of
\(SL(2,\mathbb C)\) with
\[\Lambda_{{\rm Lord}\,\nu}{}^\mu\,L_{\rm Lord}\,\sigma^{\rm LL}_\nu\,L_{\rm Lord}^\dagger=\sigma^{\rm LL}_\mu\tag{2.3}\]
(Lord (1.8)); for an infinitesimal transformation
\(\Lambda_{{\rm Lord}\,\mu}{}^\nu=\delta_\mu^\nu+\eta^{{\rm LL}\,\nu\rho}\lambda^{\rm Lord}_{\rho\mu}\)
(\(\lambda^{\rm Lord}_{\rho\mu}=-\lambda^{\rm Lord}_{\mu\rho}\)) the
solution is
\(L_{\rm Lord}=\mathsf 1_2+\tfrac14\lambda^{\rm Lord}_{\mu\nu}\sigma^{\mu\nu}_{\rm LG}\).
A Weyl spinor \(\phi\) of the kind \((\tfrac12,0)\) transforms as
\(\phi\mapsto L_{\rm Lord}\phi\), one of the kind \((0,\tfrac12)\) as
\(\chi\mapsto(L_{\rm Lord}^\dagger)^{-1}\chi\); a Dirac spinor is the
pair \((\phi,\chi)\).

\emph{Checks} (Section 9): L1-1 to L1-6 (Rust), YL-1, YL-2 (Python),
WL-1, WL-4 (Wolfram, exact): the relations (2.1), (2.2), the
self-duality, the infinitesimal solution of (2.3), and for random
matrices \(L_{\rm Lord}\) that the \(\Lambda_{\rm Lord}\) of (2.3) is a
proper orthochronous Lorentz matrix.

\hypertarget{dirac-matrices-in-the-weyl-form}{%
\subsection{2.2 Dirac matrices in the Weyl
form}\label{dirac-matrices-in-the-weyl-form}}

With the matrices (1.2),
\(\tfrac12(\gamma^\mu_{\rm LL}\gamma^\nu_{\rm LL}+\gamma^\nu_{\rm LL}\gamma^\mu_{\rm LL})=\eta^{{\rm LL}\,\mu\nu}\mathsf 1_4\),
the chirality matrix anticommutes with every \(\gamma^\mu_{\rm LL}\) and
squares to 1,
\(C_{\rm LL}\gamma^\mu_{\rm LL}C_{\rm LL}^{-1}=-(\gamma^\mu_{\rm LL})^T\),
and
\(\gamma^{\mu\nu}_{\rm LL}=\mathrm{diag}(\sigma^{\mu\nu}_{\rm LG},-(\sigma^{\mu\nu}_{\rm LG})^\dagger)\),
so that \(\tfrac12\gamma^{\mu\nu}_{\rm LL}\) generate the Lorentz
transformations of Dirac spinors (Lord (1.10) to (1.15)). The
\(4\times4\) spinor matrix is
\(\mathsf L_{\rm Lord}=\mathrm{diag}(L_{\rm Lord},(L_{\rm Lord}^\dagger)^{-1})\),
and it satisfies
\[\Lambda_{{\rm Lord}\,\nu}{}^\mu\,\mathsf L_{\rm Lord}\,\gamma^\nu_{\rm LL}\,\mathsf L_{\rm Lord}^{-1}=\gamma^\mu_{\rm LL}\tag{2.4}\]
(Lord (1.16)). The chirality matrix maps \(\gamma_{\mu\nu}\) to its
dual,
\(-\gamma^{\rm ch}_{\rm LL}\gamma_{\mu\nu\,{\rm LL}}=\gamma_{\bar\mu\bar\nu\,{\rm LL}}\)
(Lord (1.17)). \emph{Checks:} L1-7 to L1-10, YL-1, WL-1.

\hypertarget{the-equations-of-bargmann-and-wigner}{%
\subsection{2.3 The equations of Bargmann and
Wigner}\label{the-equations-of-bargmann-and-wigner}}

For spin \(s>0\) and rest mass \(m\), Lord writes the free field
equations of Bargmann and Wigner (BW) for a completely symmetric Dirac
spinor \(\psi^{\rm Lord}_{A_1\dots A_{2s}}\) of rank \(2s\) (capital
Latin letters are four-valued spinor indices):
\[\gamma^{\mu}_{{\rm LL}\,A_1}{}^{E}\,\partial_\mu\,\psi^{\rm Lord}_{EA_2\dots A_{2s}}=m\,\psi^{\rm Lord}_{A_1\dots A_{2s}},\tag{2.5}\]
and for \(m=0\) the additional condition that the chirality matrix act
on one (hence every) index as \(+1\) in his sign convention (Lord
(1.18), (1.19)): a massless field is then a completely symmetric Weyl
spinor of rank \(2s\). Spin 0 fits the scheme as the spin-1 equation
with an antisymmetric spinor. The equations are invariant under the
Lorentz group acting on every index by \(\mathsf L_{\rm Lord}\) (Lord
(1.20)).

\hypertarget{the-changes-of-representation-and-the-group-sl2mathbb-c-that-lord-gauges}{%
\subsection{\texorpdfstring{2.4 The changes of representation, and the
group \(SL(2,\mathbb C)\) that Lord
gauges}{2.4 The changes of representation, and the group SL(2,\textbackslash mathbb C) that Lord gauges}}\label{the-changes-of-representation-and-the-group-sl2mathbb-c-that-lord-gauges}}

The equations (2.5) are invariant under a second group: a change of the
representation of the Dirac matrices,
\(\gamma^\mu\mapsto S\gamma^\mu S^{-1}\), together with
\(\psi\mapsto S\psi\) on every index and \(\partial_\mu\) unchanged
(Lord (1.21)); for any invertible \(4\times4\) matrix this is the group
\(SL(4,\mathbb C)\) (up to scale). Lord keeps its subgroup that
preserves the Weyl form: the changes of representation of the Pauli
algebra,
\[\sigma_\mu\mapsto S_{\rm Lord}\,\sigma_\mu\,S_{\rm Lord}^\dagger,\qquad\phi\mapsto S_{\rm Lord}\phi,\qquad\chi\mapsto(S_{\rm Lord}^\dagger)^{-1}\chi,\tag{2.6}\]
with \(S_{\rm Lord}\) a unimodular \(2\times2\) matrix, and the bar
(1.1) taken as the definition of \(\bar\sigma_\mu\), which then changes
as
\(\bar\sigma_\mu\mapsto(S_{\rm Lord}^\dagger)^{-1}\bar\sigma_\mu S_{\rm Lord}^{-1}\)
(Lord (1.22), (1.23)). On Dirac spinors the change is the \(4\times4\)
matrix
\(\mathsf S_{\rm Lord}=\mathrm{diag}(S_{\rm Lord},(S_{\rm Lord}^\dagger)^{-1})\),
which leaves the chirality matrix unchanged. This group,
\(SL(2,\mathbb C)\), is the one that Lord makes local in his Section 4;
it is not the Lorentz group, whose parameters stay constant (Section
5.2).

\textbf{Which sets of four matrices.} Lord states that after a change
(2.6) the four matrices may be any linearly independent set of Hermitian
matrices, and that all the properties of Section 2.1 are unaltered. The
second statement holds for the sets
\(S_{\rm Lord}\sigma^{\rm LL}_\mu S_{\rm Lord}^\dagger\) (check L1-12).
The first needs a qualification, which the checks make precise: for any
four \(2\times2\) matrices, Hermitian or not,
\[\tfrac12\big(\sigma_a\bar\sigma_b+\sigma_b\bar\sigma_a\big)=N_{\rm LG}(\sigma_a,\sigma_b)\,\mathsf 1_2,\qquad N_{\rm LG}(X,Y)=\tfrac12\big(\det(X+Y)-\det X-\det Y\big),\tag{2.7}\]
because the bar is the adjugate (checks L1-13, YL-3, WL-3; exact).
\(N_{\rm LG}\) is the Minkowski product of two Hermitian matrices read
as points of space-time (the determinant of a Hermitian matrix
\(x^\mu\sigma^{\rm LL}_\mu\) is \((x^0)^2-|\vec x|^2\)). Hence (2.1),
and with it the Clifford relation of the \(\gamma\)'s, holds for a set
exactly when \(N_{\rm LG}(\sigma_a,\sigma_b)=\eta^{\rm LL}_{ab}\), which
a generic set of linearly independent Hermitian matrices violates (the
smallest deviation from \(\eta^{\rm LL}\) over twenty random sets is
\(1.07\), check L1-14). Since
\(\det(S_{\rm Lord}XS^\dagger_{\rm Lord})=\det X\), the sets that (2.6)
produces are exactly the orthonormal ones; ``any linearly independent
set'' is to be read as ``any set
\(S_{\rm Lord}\sigma^{\rm LL}_\mu S_{\rm Lord}^\dagger\)''. The general
case is not lost: a set with \(N_{\rm LG}(\sigma_a,\sigma_b)=g_{ab}\)
for another metric \(g_{ab}\) generates the Clifford algebra of
\(g_{ab}\), which is where the reference vectors of Lord's Section 4
(Section 5.6) put the metric.

\hypertarget{the-difficulty-of-the-minimal-coupling-and-lords-remedy}{%
\subsection{2.5 The difficulty of the minimal coupling, and Lord's
remedy}\label{the-difficulty-of-the-minimal-coupling-and-lords-remedy}}

If \(\partial_\mu\) is replaced by a derivative that does not commute
with itself, such as the electromagnetic \(\partial_\mu-ieA_\mu\), the
\(2s\) equations obtained from (2.5) by permuting the indices no longer
agree: for \(s>\tfrac12\) they imply an algebraic relation between the
field strength and the spin-\(s\) field that does not depend on the
strength of the coupling (Lord (1.24)), which is not acceptable. Spin
\(\tfrac12\) and massless fields are free of the defect. Lord avoids it
by applying the minimal coupling not to (2.5) but to the sum of the
permuted equations, which for spin 0 and 1 is the Kemmer form,
\[\tfrac12\big(\gamma^{\mu}_{{\rm LL}\,A}{}^{E}D_\mu\psi^{\rm Lord}_{EB}+\gamma^{\mu}_{{\rm LL}\,B}{}^{E}D_\mu\psi^{\rm Lord}_{AE}\big)=m\,\psi^{\rm Lord}_{AB}\tag{2.8}\]
(Lord (1.25)); without the coupling it is equivalent to (2.5). For
massless fields he uses the divergence form with the chirality
condition, which implies the BW equation (Lord (1.26)). The brochon is a
field of spin \(\tfrac12\), for which this difficulty does not arise.

\hypertarget{the-massless-fields-of-spin-1-and-2-lords-section-2}{%
\section{3. The massless fields of spin 1 and 2 (Lord's Section
2)}\label{the-massless-fields-of-spin-1-and-2-lords-section-2}}

\textbf{Spin 1.} For \(m=0\), \(s=\tfrac12\) the equations are those of
the massless neutrino; for \(m=0\), \(s=1\) they are Maxwell's
equations. A symmetric spinor \(\psi^{\rm Lord}_{AB}\) that satisfies
the chirality condition is a combination of the matrices
\((\gamma_{\mu\nu}+\gamma_{\bar\mu\bar\nu})C^{-1}\), with a complex
self-dual tensor \(\psi^{\rm Lord}_{\mu\nu}\) as coefficients (Lord
(2.1)); the BW equation becomes
\(\partial^\mu\psi^{\rm Lord}_{\mu\nu}=0\), and with
\(\psi^{\rm Lord}_{\mu\nu}=f_{\mu\nu}+f_{\bar\mu\bar\nu}\), \(f\) real,
it splits into \(\partial^\mu f_{\mu\nu}=0\) and
\(\partial_{[\mu}f_{\rho\sigma]}=0\) (Lord (2.2) to (2.4)). Charges
modify only the first; the second, unchanged, implies potentials,
\(f_{\mu\nu}=\partial_\mu A_\nu-\partial_\nu A_\mu\) (Lord (2.5),
(2.6)).

\textbf{Spin 2.} A completely symmetric rank-4 spinor with the chirality
condition corresponds to a complex tensor
\(\psi^{\rm Lord}_{\mu\nu\rho\sigma}\), antisymmetric in each pair,
symmetric under the exchange of the pairs, self-dual in each pair,
traceless, \(\psi^{{\rm Lord}\,\mu}{}_{\nu\rho\mu}=0\), and cyclic,
\(\psi^{\rm Lord}_{\mu\nu\rho\sigma}+\psi^{\rm Lord}_{\nu\rho\mu\sigma}+\psi^{\rm Lord}_{\rho\mu\nu\sigma}=0\)
(Lord (2.7) to (2.15)): five complex components, like a symmetric Weyl
spinor of rank 4. It is expressed by a real tensor
\(A^{\rm Lord}_{\mu\nu\rho\sigma}\) with the algebraic symmetries of the
Riemann tensor that is traceless (the symmetries of the Weyl tensor: ten
real components), and the BW equation becomes
\(\partial^\mu A^{\rm Lord}_{\mu\nu\rho\sigma}=0\) (Lord (2.18) to
(2.25)).

\textbf{Potentials of the spin-2 field.} Lord shows that, in the free
theory, \(A^{\rm Lord}_{\mu\nu\rho\sigma}\) derives from potentials
\(A^{\rm Lord}_{\rho\sigma\lambda}\), and these from ``secondary
potentials'' \(A^{\rm Lord}_{\sigma\lambda}\), which a gauge freedom of
the kind
\(A_{\rho\sigma\lambda}\mapsto A_{\rho\sigma\lambda}+\partial_\rho a_{\sigma\lambda}\)
with an antisymmetric \(a\) allows to be chosen symmetric (Lord (2.26)
to (2.36)). In terms of the symmetric \(A^{\rm Lord}_{\mu\nu}\) the
traceless condition reads
\[\Box A^{\rm Lord}_{\mu\nu}-\partial^\rho\partial_\nu A^{\rm Lord}_{\mu\rho}-\partial^\rho\partial_\mu A^{\rm Lord}_{\nu\rho}+\partial_\mu\partial_\nu A^{\rm Lord}=0,\qquad A^{\rm Lord}=\eta^{{\rm LL}\,\mu\nu}A^{\rm Lord}_{\mu\nu}\tag{3.1}\]
(Lord (2.37)): precisely the linearized empty-space equations of
Einstein, but here as the exact equations of a linear spin-2 field in
flat space-time, not as a weak-field approximation. In Einstein's theory
with \(g_{\mu\nu}=\eta_{\mu\nu}+\epsilon A_{\mu\nu}\) the same equations
hold to first order in \(\epsilon\); there it is the Riemann tensor,
\(\tfrac12\epsilon A_{\mu\nu\rho\sigma}+O(\epsilon^2)\), that is
traceless, and the BW equation \(\partial^\mu A_{\mu\nu\rho\sigma}=0\)
corresponds to the Bianchi identity.

\hypertarget{the-interacting-spin-2-field-lords-section-3}{%
\section{4. The interacting spin-2 field (Lord's Section
3)}\label{the-interacting-spin-2-field-lords-section-3}}

With interactions, a current enters the right side,
\(\partial^\mu A^{\rm Lord}_{\mu\nu\rho\sigma}=f_{\rm Lord}\,j^{\rm Lord}_{\nu\rho\sigma}\),
with a coupling constant \(f_{\rm Lord}\) (Lord (3.1)); the symmetries
of \(A\) restrict the current (Lord (3.2) to (3.5)). The potentials of
Section 3 no longer exist, because the second field equation is changed.
Lord assumes that the current has the cyclic property
\(\partial_\mu j_{\nu\rho\sigma}+\partial_\nu j_{\rho\sigma\mu}+\partial_\rho j_{\mu\nu\sigma}=0\)
(Lord (3.6)); then
\(j_{\nu\rho\sigma}=\partial_\rho j_{\sigma\nu}-\partial_\sigma j_{\rho\nu}\)
with a symmetric tensor \(j^{\rm Lord}_{\mu\nu}\) (Lord (3.7) to (3.9)),
and the tensor
\(B^{\rm Lord}_{\mu\nu\rho\sigma}=A^{\rm Lord}_{\mu\nu\rho\sigma}+f_{\rm Lord}(\eta_{\mu\rho}j_{\nu\sigma}-\eta_{\nu\rho}j_{\mu\sigma}+\eta_{\nu\sigma}j_{\mu\rho}-\eta_{\mu\sigma}j_{\nu\rho})\),
which has the symmetries of the Riemann tensor but not its
tracelessness, has potentials again (Lord (3.10) to (3.17)). The
traceless condition on \(A\) becomes
\[\Box A^{\rm Lord}_{\mu\nu}-\partial^\rho\partial_\nu A^{\rm Lord}_{\mu\rho}-\partial^\rho\partial_\mu A^{\rm Lord}_{\nu\rho}+\partial_\mu\partial_\nu A^{\rm Lord}=-f_{\rm Lord}\big(2j^{\rm Lord}_{\mu\nu}+\eta^{\rm LL}_{\mu\nu}j^{\rm Lord}\big)\tag{4.1}\]
(Lord (3.18)), the form of the linearized Einstein equations with a
source: the analogy holds with \(\eta_{\mu\nu}j-j_{\mu\nu}\)
proportional to the stress-energy tensor. Lord stresses that these are
formal similarities of a linear theory in flat space-time.

\hypertarget{the-change-of-representation-as-a-gauge-group-lords-section-4}{%
\section{5. The change of representation as a gauge group (Lord's
Section
4)}\label{the-change-of-representation-as-a-gauge-group-lords-section-4}}

This is the central section of the paper, and the one that the task
builds on. Its equations are restated here with the checks that decide
them; three of them, as printed, need a correction of a sign or a factor
(Sections 5.3, 5.4, 5.6), which the programs establish.

\hypertarget{the-local-group-and-the-dirac-matrices-as-fields}{%
\subsection{5.1 The local group and the Dirac matrices as
fields}\label{the-local-group-and-the-dirac-matrices-as-fields}}

The matrix of (2.6) is made a function of the point,
\(S_{\rm Lord}(x)\), of determinant 1, acting on Dirac spinors by
\[\mathsf S_{\rm Lord}(x)=\begin{pmatrix}S_{\rm Lord}(x)&0\\0&(S_{\rm Lord}(x)^\dagger)^{-1}\end{pmatrix}\tag{5.1}\]
(Lord (4.1)). The four Pauli-type matrices become matrix fields
\(\sigma^a_{\rm Lord}(x)\) (\(a=0,\dots,3\); Lord now uses small Latin
indices for the labels of vector components), changing as
\(\sigma^a_{\rm Lord}\mapsto S_{\rm Lord}\sigma^a_{\rm Lord}S^\dagger_{\rm Lord}\);
their bars change as in Section 2.4; and the Dirac matrices built from
them,
\[\gamma^a_{\rm Lord}(x)=\begin{pmatrix}0&\sigma^a_{\rm Lord}(x)\\\bar\sigma^a_{\rm Lord}(x)&0\end{pmatrix},\tag{5.2}\]
change as
\(\gamma^a_{\rm Lord}\mapsto\mathsf S_{\rm Lord}\gamma^a_{\rm Lord}\mathsf S_{\rm Lord}^{-1}\)
(Lord (4.2) to (4.5); check L4-1, YL-4: for
\(\sigma^a_{\rm Lord}=S_{\rm Lord}\sigma^{\rm LL}_aS_{\rm Lord}^\dagger\)
the matrices (5.2) are
\(\mathsf S_{\rm Lord}\gamma^a_{\rm LL}\mathsf S_{\rm Lord}^{-1}\) at
every point, values and derivatives). By Section 2.4 they satisfy the
Clifford relation of \(\eta^{\rm LL}\) at every point.

\hypertarget{non-commuting-derivatives-and-the-lorentz-group-as-a-group-with-constant-parameters}{%
\subsection{5.2 Non-commuting derivatives, and the Lorentz group as a
group with constant
parameters}\label{non-commuting-derivatives-and-the-lorentz-group-as-a-group-with-constant-parameters}}

The derivative of the free equations is replaced by operators
\(d^{\rm Lord}_a\) that obey the rule of Leibniz and annihilate
constants, but need not commute,
\([d^{\rm Lord}_a,d^{\rm Lord}_b]=X^{\rm Lord}_{ab}{}^c\,d^{\rm Lord}_c\)
(Lord (4.6), (4.7)); the Jacobi identity restricts \(X^{\rm Lord}\)
(Lord (4.7a)). Lord stresses that this is not yet an assumption: with
\(X^{\rm Lord}=0\) the \(d^{\rm Lord}_a\) are the partial derivatives of
Euclidean coordinates of flat space-time, and in general they could be
combinations of them; he uses only (4.6), (4.7).

A Lorentz transformation of a Dirac spinor is still given by (2.3), but
now with the point-dependent matrices \(\sigma^a_{\rm Lord}(x)\) in
place of \(\sigma^{\rm LL}_\mu\): its spinor matrix \(L_{\rm Lord}(x)\)
depends on the point although the parameters \(\Lambda_{\rm Lord}\) are
constants (Lord (4.8), (4.9)). The Lorentz group is thus
``Yang-Mills-like'' in an unusual way: through its generators, not
through its parameters. The free equations
\(\gamma^{aE}_{{\rm Lord}\,A_1}d^{\rm Lord}_a\psi^{\rm Lord}_{EA_2\dots}=m\,\psi^{\rm Lord}_{A_1\dots}\)
(Lord (4.11)) are then not invariant under the Lorentz transformations
(Lord (4.12)), because \(d^{\rm Lord}_aL_{\rm Lord}\neq0\).

\hypertarget{the-auxiliary-fields-and-the-sign-of-their-inhomogeneous-term}{%
\subsection{5.3 The auxiliary fields, and the sign of their
inhomogeneous
term}\label{the-auxiliary-fields-and-the-sign-of-their-inhomogeneous-term}}

The remedy is a generalized derivative with auxiliary fields
\(\Gamma^{\rm Lord}_a\), one matrix for each label \(a\), acting on each
spinor index:
\[D^{\rm Lord}_a\psi^{\rm Lord}_{A_1\dots A_{2s}}=d^{\rm Lord}_a\psi^{\rm Lord}_{A_1\dots A_{2s}}-\Gamma^{{\rm Lord}\,E}_{a\,A_1}\psi^{\rm Lord}_{EA_2\dots A_{2s}}-\dots-\Gamma^{{\rm Lord}\,E}_{a\,A_{2s}}\psi^{\rm Lord}_{A_1\dots E}\tag{5.3}\]
(Lord (4.13)), and \(D^{\rm Lord}_a\) of a Lorentz scalar is
\(d^{\rm Lord}_a\), which fixes its action on upper indices. The fields
are combinations of the generators,
\[\Gamma^{\rm Lord}_a=\tfrac12\,g_{\rm Lord}\,a^{\rm Lord}_{abc}\,\gamma^{bc}_{\rm Lord},\qquad\gamma^{bc}_{\rm Lord}=\tfrac12\big(\gamma^b_{\rm Lord}\gamma^c_{\rm Lord}-\gamma^c_{\rm Lord}\gamma^b_{\rm Lord}\big),\tag{5.4}\]
with a universal coupling constant \(g_{\rm Lord}\) of the auxiliary
fields \(a^{\rm Lord}_{abc}=-a^{\rm Lord}_{acb}\) to all BW fields (Lord
(4.15)). For one spinor index and a transformation
\(\psi\mapsto\mathsf T\psi\) by a matrix field \(\mathsf T(x)\), the
requirement \(D'_a(\mathsf T\psi)=\mathsf T\,D_a\psi\) and the rule of
Leibniz \(d_a(\mathsf T\psi)=\mathsf T\,d_a\psi+(d_a\mathsf T)\psi\)
give
\[\Gamma'^{\rm Lord}_a=\mathsf T\,\Gamma^{\rm Lord}_a\,\mathsf T^{-1}+(d_a\mathsf T)\,\mathsf T^{-1}.\tag{5.5}\]
Lord prints the inhomogeneous term as \(\mathsf T\,d_a\mathsf T^{-1}\)
(in his (4.14) for the Lorentz transformations and his (4.31) for the
gauge transformations). Read as the product
\(\mathsf T\,(d_a\mathsf T^{-1})\), it equals
\(-(d_a\mathsf T)\mathsf T^{-1}\): the opposite sign, which belongs to
the derivative \(d_a+\Gamma_a\). With the derivative (5.3) as he writes
it, covariance fails by a relative amount of at least \(0.16\) at the
points tested (check L4-3; YL-4b; exactly, WL-5), while (5.5) holds to
the rounding (L4-2, YL-4, WL-5). The two forms differ only by the name
\(\Gamma\leftrightarrow-\Gamma\), so no physics depends on the choice;
this manuscript keeps Lord's derivative (5.3) and the law (5.5).

\textbf{The coefficients and the trace.} Because
\(\mathrm{tr}(\gamma_{bc}\gamma^{de})=-4(\delta_b^d\delta_c^e-\delta_b^e\delta_c^d)\)
(for \(b\neq c\)), the coefficients of (5.4) are
\(a^{\rm Lord}_{abc}=-\frac1{g_{\rm Lord}}\mathrm{tr}_{\rm N}(\gamma^{\rm Lord}_{bc}\Gamma^{\rm Lord}_a)\),
and under a gauge transformation (5.5) they change by
\[\delta a^{\rm Lord}_{abc}=-\frac1{g_{\rm Lord}}\,\mathrm{tr}_{\rm N}\big(\gamma'^{\rm Lord}_{bc}\,(d_a\mathsf T)\mathsf T^{-1}\big),\tag{5.6}\]
with \(\gamma'\) the transformed matrices (checks L4-4, YL-5). Lord's
(4.16) and (4.32) have this form with ``tr'' for
\(\mathrm{tr}_{\rm N}\), the trace normalized to 1 on the unit matrix,
and with his inhomogeneous term: as printed, the change is exactly the
negative of (5.6) (the printed change minus the true one is \(2.000000\)
times the true one, check L4-5). So Lord's (4.15), (4.16) and (4.31),
(4.32) are consistent with each other when tr is the normalized trace;
their common sign is that of the derivative \(d_a+\Gamma_a\).

\hypertarget{the-connection-of-the-vectors-and-a-factor-2}{%
\subsection{5.4 The connection of the vectors, and a factor
2}\label{the-connection-of-the-vectors-and-a-factor-2}}

For vectors Lord defines
\(D^{\rm Lord}_av_b=d^{\rm Lord}_av_b-C^{\rm Lord}_{ba}{}^e v_e\) with a
tensorial \(C^{\rm Lord}\) (Lord (4.17), (4.18)), chosen so that the
Dirac matrices are covariantly constant,
\[D^{\rm Lord}_a\gamma^{\rm Lord}_b=d^{\rm Lord}_a\gamma^{\rm Lord}_b-[\Gamma^{\rm Lord}_a,\gamma^{\rm Lord}_b]-C^{\rm Lord}_{ba}{}^e\gamma^{\rm Lord}_e=0\tag{5.7}\]
(Lord (4.19)). This is possible because both \(d_a\gamma_b\) and
\([\Gamma_a,\gamma_b]\) are combinations of the \(\gamma_e\) (check
L4-6a). From
\([\gamma^{cd},\gamma_b]=2(\gamma^c\delta^d_b-\gamma^d\delta^c_b)\)
follows
\([\Gamma^{\rm Lord}_a,\gamma^{\rm Lord}_b]=-2g_{\rm Lord}\,a^{\rm Lord}_{ab}{}^e\gamma^{\rm Lord}_e\),
and therefore
\[C^{\rm Lord}_{ba}{}^e=2\,g_{\rm Lord}\,a^{\rm Lord}_{ab}{}^e+\mathrm{tr}_{\rm N}\big[(d^{\rm Lord}_a\gamma^{\rm Lord}_b)\gamma^e_{\rm Lord}\big]\tag{5.8}\]
(checks L4-6b, YL-6, and exactly in the special gauge, WL-6). Lord
prints (4.20) as
\(C_{ba}{}^e=g\,a_{ab}{}^e+\tfrac12\mathrm{tr}[(d_a\gamma_b)\gamma^e]\).
With tr read as \(\mathrm{tr}_{\rm N}\), as Section 5.3 requires, this
is exactly half of (5.8) (check L4-6c); with the full trace it is not
even proportional (the largest coefficient deviates by \(0.56\) of the
largest one, check L4-6d). No normalization of the trace can repair
(4.20), because its term \(g\,a_{ab}{}^e\) has half the coefficient
\(2g\) that (5.7) requires. The properties that Lord derives from (5.7)
hold for (5.8): \(C^{\rm Lord}\) is a tensor under the Lorentz
transformations, it is antisymmetric,
\(C^{\rm Lord}_{bae}=-C^{\rm Lord}_{eab}\) (Lord (4.24); check L4-7), it
leaves the Levi-Civita symbol constant, \(D_e\epsilon_{abcd}=0\) (Lord
(4.25); check L4-8), and it is invariant under the gauge transformations
of Section 5.5. Section 5.6 shows the consequence of the factor for the
curvature.

With (5.7) and the constancy of \(\epsilon\), the generalized BW
equations
\(\gamma^{aE}_{{\rm Lord}\,A_1}D^{\rm Lord}_a\psi^{\rm Lord}_{EA_2\dots}=m\psi^{\rm Lord}_{A_1\dots}\)
pass to the tensor form of Section 3 by replacing \(\partial_a\) with
\(D^{\rm Lord}_a\) (Lord (4.26) to (4.29)); but the potentials of the
spin-2 field can then no longer be introduced as in Section 3, because
the \(D^{\rm Lord}_a\) do not commute. Lord resolves this in his Section
5 by identifying the auxiliary fields \(a^{\rm Lord}_{abc}\) themselves
as the potentials of the spin-2 field.

\hypertarget{the-gauge-invariance}{%
\subsection{5.5 The gauge invariance}\label{the-gauge-invariance}}

The equations generalized by \(d_a\to D^{\rm Lord}_a\) are invariant
under the local changes of representation:
\(\psi\mapsto\mathsf S_{\rm Lord}\psi\) on every index,
\(\gamma^a_{\rm Lord}\mapsto\mathsf S_{\rm Lord}\gamma^a_{\rm Lord}\mathsf S_{\rm Lord}^{-1}\),
the chirality matrix and \(d^{\rm Lord}_a\) unchanged, for any
unimodular field \(S_{\rm Lord}(x)\), provided that the auxiliary fields
change by (5.5) and (5.6) (Lord (4.30) to (4.32), with the sign of
Section 5.3). The \(C^{\rm Lord}_{abc}\) do not change. This group,
\(SL(2,\mathbb C)\) made local, is \textbf{Lord's gauge group}: the
local group of the changes of representation of the Clifford algebra of
the Dirac matrices that keep its Weyl form.

\hypertarget{reference-vectors-the-metric-the-affine-connection-and-the-curvature}{%
\subsection{5.6 Reference vectors, the metric, the affine connection and
the
curvature}\label{reference-vectors-the-metric-the-affine-connection-and-the-curvature}}

Lord introduces at every point four real, linearly independent
\textbf{reference vectors} \(h^{\rm Lord}_a{}^\mu\) (a tetrad; the Greek
index labels the four vectors), with
\(g^{\mu\nu}_{\rm Lord}=h^{\rm Lord}_a{}^\mu h^{{\rm Lord}\,a\nu}\) and
its inverse \(g^{\rm Lord}_{\mu\nu}\), which raise and lower Greek
indices, so that \(v_a=h^{\rm Lord}_a{}^\mu v_\mu\) and
\(\eta^{\rm LL}_{ab}=h^{\rm Lord}_a{}^\mu h^{\rm Lord}_{b\mu}\) (Lord
(4.33) to (4.36)). The derivative is extended to them by requiring
\(D_ah_b{}^\rho=d_ah_b{}^\rho-C_{ba}{}^\rho+\Gamma_{ba}{}^\rho=0\),
which defines the coefficients \(\Gamma^{\rm Lord}_{\mu\nu}{}^\rho\)
(Lord (4.37), (4.38)). The formalism is invariant under the
transformations \(h_a{}^\mu\mapsto G^\mu{}_\nu h_a{}^\nu\) of the
reference vectors (a local \(SL(4,\mathbb R)\)); restricting them by
\(d_\rho G^\mu{}_\nu-d_\nu G^\mu{}_\rho=0\) allows the \(d_\mu\) to be
chosen commuting, so that they are partial derivatives of coordinates,
and the restricted transformations are then the general coordinate
transformations, under which \(\Gamma^{\rm Lord}_{\mu\nu}{}^\rho\)
transforms as an affine connection (Lord (4.39) to (4.42)). Its
symmetric part is the Christoffel symbol plus terms of the torsion
\(T^{\rm Lord}_{\mu\nu}{}^\rho\), its antisymmetric part is
\(\tfrac12T^{\rm Lord}_{\mu\nu}{}^\rho\) (Lord (4.43), (4.44)). The
commutator of two derivatives gives the curvature: on a vector,
\[[D_\mu,D_\nu]v_\rho=R^{\rm Lord}_{\mu\nu\rho}{}^\sigma v_\sigma+T^{\rm Lord}_{\mu\nu}{}^\sigma D_\sigma v_\rho,\qquad
R^{\rm Lord}_{\mu\nu\rho}{}^\sigma=\partial_\nu\Gamma_{\rho\mu}{}^\sigma-\partial_\mu\Gamma_{\rho\nu}{}^\sigma+\Gamma_{\rho\mu}{}^\lambda\Gamma_{\lambda\nu}{}^\sigma-\Gamma_{\rho\nu}{}^\lambda\Gamma_{\lambda\mu}{}^\sigma\tag{5.9}\]
(Lord (4.45), (4.46)), and on a spinor with one lower index, when the
torsion vanishes,
\[[D^{\rm Lord}_\mu,D^{\rm Lord}_\nu]\psi=\big(\partial_\nu\Gamma^{\rm Lord}_\mu-\partial_\mu\Gamma^{\rm Lord}_\nu+[\Gamma^{\rm Lord}_\mu,\Gamma^{\rm Lord}_\nu]\big)\psi\tag{5.10}\]
(checks L4-9, YL-7, with exact second derivatives in Rust and nested
differences in Python). Lord writes the bracket of (5.10) as
\(\tfrac12R^{\rm Lord}_{\mu\nu\alpha\beta}\gamma^{\alpha\beta}\) (his
(4.47), (4.48)). The checks show that the correct coefficient is
\(\tfrac14\): if the spinor connection keeps the Dirac matrices
constant, \([\Gamma_\mu,\gamma_b]=-W_{\mu\,b}{}^e\gamma_e\) for the
connection \(W_\mu\) of the vectors, then
\(\Gamma_\mu=\tfrac14W_{\mu\,ab}\gamma^{ab}\), the map
\(W\mapsto\tfrac14W_{ab}\gamma^{ab}\) is a homomorphism of Lie algebras
(exact, WL-7), and the bracket of (5.10) is
\[\partial_\nu\Gamma_\mu-\partial_\mu\Gamma_\nu+[\Gamma_\mu,\Gamma_\nu]=\tfrac14R_{\mu\nu\alpha\beta}\,\gamma^{\alpha\beta}\tag{5.11}\]
to all orders (checks L4-10a, L4-10b, YL-7b): Lord's \(\tfrac12\) gives
twice the spinor curvature (the relative deviation is \(1.000000\),
check L4-10c). The two printed factors, \(\tfrac12\) in (4.20) and
\(\tfrac12\) in (4.48), compensate each other to first order: with
\(C=\tfrac12W\) and \(\tfrac12R(C)\gamma\) the curvature is reproduced
up to terms of second order in the connection, which fall by a factor
\(96.0\) when the connection falls by 10 (check L4-10d). The non-linear
theory of Lord's Section 5 needs the exact relations (5.8) and (5.11).
The Jacobi identity of the \(D_\mu\) gives the Bianchi identities with
torsion (Lord (4.49) to (4.51)). With
\(\Gamma^{\rm Lord}_{\mu\nu}{}^\rho\) an affine connection, the
\(g^{\rm Lord}_{\mu\nu}\) are the components of the metric, which
becomes a physical field: the coupling (5.3), with the requirement that
the \(d_a\) be combinations of the \(\partial_\mu\), makes space-time
curved.

\hypertarget{the-non-linear-spin-2-theory-lords-section-5}{%
\section{6. The non-linear spin-2 theory (Lord's Section
5)}\label{the-non-linear-spin-2-theory-lords-section-5}}

\textbf{The flat limit.} The generalized BW theory
\(\gamma^{\mu E}_{A_1}D_\mu\psi_{EA_2\dots}=m\psi_{A_1\dots}\) (Lord
(5.1)) returns to the flat theory of Sections 2 to 4 as
\(g_{\rm Lord}\to0\). To follow the limit, Lord fixes a special gauge in
which the Dirac matrices \(\gamma_a\) are constant; then (5.7) gives
\(C_{bae}=g\,a_{abe}\) by his (4.20), and \(2g\,a_{abe}\) by (5.8) (Lord
(5.2); WL-6). He accompanies a change of the coupling by a change of the
reference vectors, \(h_a{}^\rho=\delta_a^\rho+g\,a_a{}^\rho\), and of
the metric, \(g_{\mu\nu}=\eta_{\mu\nu}+g\,A_{\mu\nu}\) with
\(A_{\mu\nu}\) finite in the limit (Lord (5.3) to (5.7)); such a change
is not a coordinate transformation, but a change of the shape of
space-time. In the limit the metric is flat, the Dirac matrices are
constant and every \(D_\mu\) becomes \(\partial_\mu\).

\textbf{The auxiliary fields are the potentials of the spin-2 field.} In
curved space-time the spin-2 equation gives only
\(D_{(\mu}A_{\nu\rho)\sigma\lambda}=0\) (Lord (5.8)), which does not
define potentials. Lord resolves this by identifying the auxiliary
fields \(a_{\mu\nu\rho}\) themselves as the potentials: the Riemann
tensor has the cyclic symmetry and satisfies (5.8) provided that the
torsion vanishes, which also removes the difficulty of the potentials of
the spin-1 field (Lord (5.9), (5.10)). With
\(\Gamma_{\mu\nu\rho}=g\,B_{\mu\nu\rho}\), the \(B_{\mu\nu\rho}\) are
the Christoffel symbols of \(A_{\mu\nu}\) (Lord (5.11), (5.12)),
\(A_{\mu\nu\rho\sigma}=\tfrac2gR_{\mu\nu\rho\sigma}\) stays finite (Lord
(5.13), (5.14)), and in the limit the ``basic Yang-Mills fields''
\(2a_{\nu\mu\rho}\) differ from the potentials of the linear spin-2
field only by a gauge transformation of Section 3 (Lord (5.15), (5.16)).
The factor 2 in \(2a\) agrees with (5.8), in which the connection of the
vectors is \(2g\,a\).

\textbf{Einstein's equations.} The remaining condition, that the spin-2
tensor be traceless, becomes \(R_{\mu\nu}=0\), Einstein's empty-space
equations (Lord (5.17)). Lord remarks that they arise here from the
chirality condition on the massless spin-2 spinor (Lord (5.18)), as the
left-handedness of the neutrino arises from the same condition for spin
\(\tfrac12\). With interaction,
\(D^\mu A_{\mu\nu\rho\sigma}=g\,j_{\nu\rho\sigma}\) (the coupling of
Section 4 identified with \(g_{\rm Lord}\)), and
\(B_{\mu\nu\rho\sigma}=\tfrac2gR_{\mu\nu\rho\sigma}\) (Lord (5.19) to
(5.21)); the contraction gives
\[R_{\mu\nu}-\tfrac12g_{\mu\nu}R=-g_{\rm Lord}^2\big(j_{\mu\nu}-g_{\mu\nu}j\big)\tag{6.1}\]
(Lord (5.22) to (5.24)), Einstein's equation with a source, once
\(j_{\mu\nu}-g_{\mu\nu}j\) is identified with a multiple of the
stress-energy tensor of matter; its divergence vanishes by the Bianchi
identity (Lord (5.25)). The spin-2 tensor is then \(\tfrac2g\) times the
Weyl tensor (Lord (5.26) to (5.30)).

\textbf{The tetrad and the local Lorentz group.} The \(h_a{}^\mu\) are
the usual vierbein fields, and the formalism is invariant under
point-dependent Lorentz rotations of the vierbein if the \(a_{abc}\) are
regarded as auxiliary fields of that group too (the \(C_{abc}\) are then
not tensors of it). Lord contrasts his approach with Sciama's, who
varies a Lagrangian \(hR\) with respect to the vierbein and the
connection independently and finds torsion in the presence of spin; Lord
abandons the Lagrangian approach, relies on the BW equations, identifies
the \(a_{\mu a}{}^b\) as the gravitational field, and makes the
vanishing of the torsion an essential part of the theory.

\hypertarget{remarks-on-the-equations-as-printed}{%
\section{7. Remarks on the equations as
printed}\label{remarks-on-the-equations-as-printed}}

The following table collects the equations of Lord's Section 4 whose
printed form does not follow from his definitions (5.3), (5.4) and
(5.7). None of them changes his conclusions, because each is a
convention or a factor that the later sections absorb (the sign of
\(\Gamma\), the normalization of \(a\)); but the non-linear theory needs
the corrected relations, and the coupling of the brochon in Part III is
built on them.

\sptkeptpenalty=10000
\begin{longtable}[]{@{}
  >{\raggedright\arraybackslash}p{(\columnwidth - 6\tabcolsep) * \real{0.2500}}
  >{\raggedright\arraybackslash}p{(\columnwidth - 6\tabcolsep) * \real{0.2500}}
  >{\raggedright\arraybackslash}p{(\columnwidth - 6\tabcolsep) * \real{0.2500}}
  >{\raggedright\arraybackslash}p{(\columnwidth - 6\tabcolsep) * \real{0.2500}}@{}}
\toprule\noalign{}
\begin{minipage}[b]{\linewidth}\raggedright
Lord
\end{minipage} & \begin{minipage}[b]{\linewidth}\raggedright
as printed
\end{minipage} & \begin{minipage}[b]{\linewidth}\raggedright
as it follows from his definitions
\end{minipage} & \begin{minipage}[b]{\linewidth}\raggedright
checks
\end{minipage} \\
\midrule\noalign{}
\endhead
\bottomrule\noalign{}
\endlastfoot
(4.14), (4.31) & inhomogeneous term \(\mathsf S\,d_a\mathsf S^{-1}\) &
\((d_a\mathsf S)\mathsf S^{-1}=-\mathsf S\,(d_a\mathsf S^{-1})\) for
\(D_a=d_a-\Gamma_a\), (5.5) & L4-2, L4-3, YL-4, WL-5 \\*
(4.16), (4.32) &
\(-\frac1g\mathrm{tr}[\gamma_{bc}\mathsf Sd_a\mathsf S^{-1}]\) &
\(-\frac1g\mathrm{tr}_{\rm N}[\gamma'_{bc}(d_a\mathsf S)\mathsf S^{-1}]\),
(5.6): tr is the normalized trace, and the sign follows (5.5) & L4-4,
L4-5, YL-5 \\*
(4.20), (5.2) &
\(C_{ba}{}^e=g\,a_{ab}{}^e+\tfrac12\mathrm{tr}[(d_a\gamma_b)\gamma^e]\)
&
\(C_{ba}{}^e=2g\,a_{ab}{}^e+\mathrm{tr}_{\rm N}[(d_a\gamma_b)\gamma^e]\),
(5.8): twice the printed form & L4-6a to L4-6d, YL-6, WL-6 \\*
(4.47), (4.48) & \(\tfrac12R_{\mu\nu\alpha\beta}\gamma^{\alpha\beta}\) &
\(\tfrac14R_{\mu\nu\alpha\beta}\gamma^{\alpha\beta}\), (5.11) & L4-10a
to L4-10d, YL-7b, YL-7c, WL-7 \\*
Section 1 & the four Hermitian matrices ``can be any linearly
independent set'' & any set \(S\sigma_\mu S^\dagger\) (\(\det S=1\)):
exactly the sets with \(N_{\rm LG}(\sigma_a,\sigma_b)=\eta_{ab}\), (2.7)
& L1-12 to L1-14, YL-3, WL-3 \\
\end{longtable}

The index structure of the homogeneous term of Lord (4.16), which
transforms only two of the three indices of \(a_{efg}\), is presumably a
misprint as well (a third factor \(\Lambda^{-1}\) belongs to the third
index); no check depends on it.

\hypertarget{what-part-i-establishes-for-the-task}{%
\section{8. What Part I establishes for the
task}\label{what-part-i-establishes-for-the-task}}

\begin{enumerate}
\def\labelenumi{\arabic{enumi}.}
\tightlist
\item
  Lord's gauge group is the group \(SL(2,\mathbb C)\) of the changes of
  representation of the Pauli algebra, made local:
  \(\gamma^a(x)\mapsto\mathsf S(x)\gamma^a(x)\mathsf S(x)^{-1}\) with
  \(\mathsf S=\mathrm{diag}(S,(S^\dagger)^{-1})\), \(\det S=1\). It acts
  on the Dirac matrices as fields, not on the coordinates; the Lorentz
  group keeps constant parameters.
\item
  The covariant derivative of a spinor is \(D_a=d_a-\Gamma_a\) with
  \(\Gamma_a\) in the Lorentz algebra and the law (5.5).
\item
  If the Dirac matrices are covariantly constant, the spinor connection
  is \(\Gamma_\mu=\tfrac14W_{\mu\,ab}\gamma^{ab}\), with \(W_\mu\) the
  connection of the vectors, and the spinor curvature is
  \(\tfrac14R_{\mu\nu ab}\gamma^{ab}\) exactly.
\item
  With reference vectors (a tetrad) and vanishing torsion, the auxiliary
  fields are the potentials of a self-interacting spin-2 field whose
  equations are Einstein's; the coupling to matter enters through (6.1).
\end{enumerate}

These are the facts that Part III uses to couple the brochon to
gravitation.

\hypertarget{provenance-of-part-i}{%
\section{9. Provenance of Part I}\label{provenance-of-part-i}}

\textbf{The input.} E. A. Lord, ``General relativity from gauge
invariance'', Proc. Camb. Phil. Soc. 69, 423 (1971): twenty pages (the
journal pages 423 to 442), a scan with recognized text; all pages were
read. The file is not in the repository; its SHA-256 checksum:

\begin{verbatim}
c928b11e5721bb37093f9b19b19e8e9d4a743fcd3f8f4779e804845fcaff2275  General relativity from gauge invariance.pdf
\end{verbatim}

\textbf{The programs.} Three programs check Part I and Part II, each
independently of the others:

\begin{itemize}
\tightlist
\item
  the Rust program \texttt{lg\_lord} (\texttt{rust/lglord/src/main.rs},
  with the library \texttt{rust/lgcore}, whose module \texttt{lord.rs}
  builds Lord's matrices from his definitions): ALL PASS (45 of 45). Its
  local checks use fields of the four coordinates built from smooth
  pseudo-random functions, with exact first and second derivatives by
  forward-mode automatic differentiation;
\item
  the Python program \texttt{python/lg\_lord.py} (with
  \texttt{python/lgcommon.py}): ALL PASS (16 of 16); it writes every
  matrix again from the definitions and differentiates by central finite
  differences of eighth order;
\item
  the Wolfram script \texttt{wolfram/source/lg\_lord.wls}: ALL PASS (7
  of 7), in exact arithmetic (symbolic matrices).
\end{itemize}

A finding about an equation as printed is a check that passes when the
printed form fails by the stated amount; its numbers are in
\texttt{data/lg\_lord\_findings.csv}. To reproduce, from the folder
\texttt{Revision\_Lord\_Gravity}:

\begin{Shaded}
\begin{Highlighting}[]
\BuiltInTok{cd} \StringTok{"/Users/nsh/Developer/github/spacetimes/Revision\_Lord\_Gravity"}
\KeywordTok{(}\BuiltInTok{cd}\NormalTok{ rust }\KeywordTok{\&\&} \ExtensionTok{cargo}\NormalTok{ build }\AttributeTok{{-}{-}release} \KeywordTok{\&\&} \ExtensionTok{./target/release/lg\_lord}\KeywordTok{)}   \CommentTok{\# data/hck\_lg\_lord\_checks.txt, lg\_lord\_findings.csv}
\ExtensionTok{python3}\NormalTok{ python/lg\_lord.py                                         }\CommentTok{\# data/hck\_lg\_lord\_checks\_python.txt}
\ExtensionTok{wolframscript} \AttributeTok{{-}file}\NormalTok{ wolfram/source/lg\_lord.wls data/hck\_lg\_lord\_checks\_wolfram.txt}
\ExtensionTok{python3}\NormalTok{ code/scripts/write\_lg\_lord\_fragments.py                  }\CommentTok{\# the numbers of this manuscript}
\end{Highlighting}
\end{Shaded}

\textbf{The report of the Rust program}
(\texttt{data/hck\_lg\_lord\_checks.txt}):

\begin{verbatim}
PROJECT spacetimes, Revision_Lord_Gravity: lg_lord -- the checks of the equations of E. A. Lord, General relativity from gauge invariance (1971), Part I of lord_gravity.md, and of the association of the SL(2,C) of the project with his gauge group, Part II
Each line: PASS or FAIL, the identifier of the check, the measured deviation and its tolerance (or, for a
statement that is true or false, the facts that decide it), and the statement that is checked.

PASS L1-1    deviation 0.000e0 (tolerance 1.0e-15): (1.2): sigma-bar^mu = -epsilon (sigma^mu)^T epsilon is (1, -sigma_1, -sigma_2, -sigma_3) for the standard set (1.1)
PASS L1-2    deviation 0.000e0 (tolerance 1.0e-15): (1.3): 1/2 (sigma^mu sigma-bar^nu + sigma^nu sigma-bar^mu) = eta^{mu nu} 1 (16 pairs)
PASS L1-3    deviation 0.000e0 (tolerance 1.0e-14): (1.5): [sigma^{mu nu}, sigma^{rho sigma}] = 2 (eta^{nu rho} sigma^{mu sigma} - eta^{mu rho} sigma^{nu sigma} + eta^{mu sigma} sigma^{nu rho} - eta^{nu sigma} sigma^{mu rho}) (256 index combinations)
PASS L1-4    deviation 0.000e0 (tolerance 1.0e-15): (1.6), (1.7): sigma_{mu nu} is self-dual, sigma_{mu nu} = 1/2 i eps_{mu nu rho sigma} sigma^{rho sigma} with eps_0123 = +1
PASS L1-5    deviation 2.220e-16 (tolerance 1.0e-15): (1.8), infinitesimal: with Lambda_mu^nu = delta_mu^nu + eta^{nu rho} lambda_{rho mu}, L = 1 + 1/4 lambda_{mu nu} sigma^{mu nu} satisfies (1.8) to first order (20 random lambda): eta^{mu rho} lambda_{rho nu} sigma^nu + X sigma^mu + sigma^mu X^dag = 0
PASS L1-6a   deviation 9.770e-15 (tolerance 1.0e-12): (1.8): for 50 random L of SL(2,C), the Lambda defined by Lambda_nu^mu L sigma^nu L^dag = sigma^mu satisfies Lambda eta Lambda^T = eta
PASS L1-6b   (1.8): ... and has det Lambda = 1 and Lambda_0^0 >= 1 (proper orthochronous): L -> Lambda is the covering map of the proper orthochronous Lorentz group [50 random L]
PASS L1-6c   deviation 9.115e-15 (tolerance 1.0e-12): (1.16): Lambda_nu^mu L4 gamma^nu L4^{-1} = gamma^mu with L4 = diag(L, (L^dag)^{-1}), the same Lambda (50 random L)
PASS L1-7a   deviation 0.000e0 (tolerance 1.0e-15): (1.10): 1/2 (gamma^mu gamma^nu + gamma^nu gamma^mu) = eta^{mu nu} for the Weyl representation (1.9)
PASS L1-7b   deviation 0.000e0 (tolerance 1.0e-15): (1.11): gamma^mu gamma^5 + gamma^5 gamma^mu = 0
PASS L1-7c   deviation 0.000e0 (tolerance 1.0e-15): (1.12): (gamma^5)^2 = 1
PASS L1-7d   deviation 0.000e0 (tolerance 1.0e-14): (1.13): gamma^5 = -i gamma^0 gamma^1 gamma^2 gamma^3 = -(1/24) i eps_{mu nu rho sigma} gamma^mu gamma^nu gamma^rho gamma^sigma (both forms)
PASS L1-8    deviation 0.000e0 (tolerance 1.0e-15): (1.14): C gamma^mu C^{-1} = -(gamma^mu)^T with C = diag(epsilon, -epsilon)
PASS L1-9    deviation 0.000e0 (tolerance 1.0e-15): (1.15): gamma^{mu nu} = diag(sigma^{mu nu}, -(sigma^{mu nu})^dag) in the Weyl representation
PASS L1-10   deviation 0.000e0 (tolerance 1.0e-15): (1.17): gamma_5 gamma_{mu nu} = gamma_{mu nu} gamma_5 = gamma_{mu-bar nu-bar} (the dual (1.6))
PASS L1-11a  deviation 8.113e-15 (tolerance 1.0e-13): property of epsilon used for (1.23): epsilon S = (S^T)^{-1} epsilon for every unimodular S (50 random S)
PASS L1-11b  deviation 5.329e-14 (tolerance 1.0e-12): (1.22), (1.23): for sigma^mu -> S sigma^mu S^dag, the sigma-bar defined by (1.2) becomes (S^dag)^{-1} sigma-bar^mu S^{-1} (50 random S)
PASS L1-12   deviation 5.440e-15 (tolerance 1.0e-12): Section 1, 'all the properties ... are unaltered': every set S sigma^mu S^dag (S unimodular) satisfies (1.3), (1.7) and (1.10) (20 random S)
PASS L1-13   deviation 5.551e-16 (tolerance 1.0e-14): for ANY four Hermitian matrices, 1/2 (sigma^a sigma-bar^b + sigma^b sigma-bar^a) = N(sigma^a, sigma^b) 1, with N the polarized determinant (the Minkowski product of the matrices read as points of space-time) (20 random sets)
PASS L1-14   hence (1.3) and (1.10) hold for a set exactly when N(sigma^a, sigma^b) = eta^{ab}: a generic linearly independent Hermitian set violates them, and Lord's 'any linearly independent set' must be read as 'any set S sigma^mu S^dag' [smallest deviation from eta over 20 random sets: 1.073]
PASS L4-1    deviation 6.661e-15 (tolerance 1.0e-12): (4.2) to (4.5): for a field S(x) of SL(2,C), the Dirac matrices gamma^a(x) built by (4.3), (4.5) from sigma^a(x) = S sigma^a S^dag equal S4 gamma^a S4^{-1}, S4 = diag(S, (S^dag)^{-1}) (4.1): values and first derivatives at 12 random points
PASS L4-2    deviation 7.183e-16 (tolerance 1.0e-13): (4.13), (4.31) corrected: D_a = d_a - Gamma_a is covariant, D'_a (T psi) = T D_a psi, when Gamma'_a = T Gamma_a T^{-1} + (d_a T) T^{-1} (relative deviation, 12 points, 4 directions)
PASS L4-3    (4.14), (4.31) as printed: with the inhomogeneous term T d_a T^{-1} read as the product T (d_a T^{-1}) = -(d_a T) T^{-1}, D_a = d_a - Gamma_a is NOT covariant; the printed sign belongs to the derivative d_a + Gamma_a [smallest relative deviation of the printed law: 0.156]
PASS L4-4    deviation 7.427e-15 (tolerance 1.0e-11): (4.15), (4.16), (4.32) corrected: Gamma_a = 1/2 g a_abc gamma^{bc} with a_abc = -(1/(4g)) tr(gamma_bc Gamma_a) (reconstruction), and under the gauge transformation delta a_abc = -(1/g) tr_N[gamma'_bc (d_a T) T^{-1}], with tr_N the trace normalized to tr_N 1 = 1 and gamma' the transformed matrices
PASS L4-5    (4.16), (4.32) as printed, with tr = tr_N: the inhomogeneous term -(1/g) tr[gamma_bc T d_a T^{-1}] has the opposite sign of the correct one (it is the law of the derivative d_a + Gamma_a) [smallest |printed - true| / |true| over all coefficients: 2.0000 (2 for an opposite sign)]
PASS L4-6a   deviation 1.427e-14 (tolerance 1.0e-12): (4.19): d_a gamma_b - [Gamma_a, gamma_b] is a linear combination of the gamma_e with real coefficients C_ba^e, as Lord states (residual and imaginary parts)
PASS L4-6b   deviation 3.997e-15 (tolerance 1.0e-12): (4.19) solved: C_ba^e = 2 g a_ab^e + tr_N[(d_a gamma_b) gamma^e] (tr_N = 1/4 tr), for Gamma_a = 1/2 g a_abc gamma^{bc} (4.15)
PASS L4-6c   deviation 1.998e-15 (tolerance 1.0e-12): (4.20) as printed, read with tr = tr_N, g a_ab^e + 1/2 tr_N[(d_a gamma_b) gamma^e], is exactly HALF of the C_ba^e that (4.19) requires
PASS L4-6d   (4.20) as printed, read with tr the full trace, is not C_ba^e either; since its term g a_ab^e carries half the coefficient 2 g that (4.19) requires, no normalization of tr makes (4.20) equal to C, and only tr_N makes it proportional (L4-6c) [smallest, over the 12 points, of the largest deviation of a coefficient relative to the largest coefficient: 0.5649]
PASS L4-7    deviation 8.438e-15 (tolerance 1.0e-12): (4.24): C_bae = -C_eab, for the C that solves (4.19)
PASS L4-8    deviation 3.580e-15 (tolerance 1.0e-11): (4.25): D_e eps_abcd = C_ae^f eps_fbcd + C_be^f eps_afcd + C_ce^f eps_abfd + C_de^f eps_abcf = 0
PASS L4-9    deviation 9.486e-16 (tolerance 1.0e-12): (4.47), (4.48): for D_mu = d_mu - Gamma_mu (coordinate derivatives, no torsion), [D_mu, D_nu] psi = (d_nu Gamma_mu - d_mu Gamma_nu + [Gamma_mu, Gamma_nu]) psi (12 points, exact second derivatives)
PASS L4-10a  deviation 2.220e-16 (tolerance 1.0e-13): the spinor image of a vector connection: Gamma_mu = 1/4 omega_{mu ab} gamma^{ab} satisfies [Gamma_mu, gamma_b] = -(W_mu)_b^e gamma_e, that is D_mu gamma_b = 0 with constant gamma (Lord's special gauge (5.2))
PASS L4-10b  deviation 1.144e-16 (tolerance 1.0e-12): then the spinor curvature is d_nu Gamma_mu - d_mu Gamma_nu + [Gamma_mu, Gamma_nu] = 1/4 R_{mu nu alpha beta} gamma^{alpha beta}, with R the curvature (4.46) of the vector connection: exactly, to all orders
PASS L4-10c  (4.47), (4.48) as printed, 1/2 R_{mu nu alpha beta} gamma^{alpha beta}, with R the curvature of the same (compatible) vector connection: twice the spinor curvature [smallest relative deviation of the printed 1/2: 1.0000 (1 for a factor 2)]
PASS L4-10d  with Lord's printed factors together, C = W/2 from (4.20) and 1/2 R(C) gamma^{ab} in (4.48), the spinor curvature is reproduced only to first order in the connection: the deviation is of second order (it falls by 100 when the connection falls by 10) [deviation at amplitude 0.4: 1.442e-1; at 0.04: 1.502e-3; ratio 96.00]
PASS P-1     deviation 0.000e0 (tolerance 0.0e0): Lord's sigma^mu (1.1) and sigma-bar^mu (1.2) are the project's sigma^LL_p and sigma-bar_p (the Pauli matrices of the chiral basis)
PASS P-2     deviation 0.000e0 (tolerance 0.0e0): Lord's epsilon (1.1) is the project's E_LL
PASS P-3     deviation 0.000e0 (tolerance 0.0e0): Lord's gamma^mu of the Weyl representation (1.9) are the project's gamma^p_LL of the chiral basis, and his C (1.14) is C_LL
PASS P-4     deviation 0.000e0 (tolerance 1.0e-15): Lord's gamma^5 (1.9) is minus the project's chirality matrix gamma^ch_LL = i gamma^0 gamma^1 gamma^2 gamma^3
PASS P-5     deviation 5.551e-16 (tolerance 1.0e-13): Lord's 4x4 change of representation (4.1), diag(S, (S^dag)^{-1}) with S = M_LL(g-tilde) = (M_LL(g)^dag)^{-1}, is the project's spinor matrix S_LL(g) = diag(M_LL(g-tilde), M_LL(g)) of the unit bi-quaternion g (50 random g)
PASS P-6a    deviation 0.000e0 (tolerance 1.0e-14): the project's conjugation of bi-quaternions is Lord's sigma-bar: M_LL(b-bar) = -epsilon M_LL(b)^T epsilon for every bi-quaternion b (50 random b)
PASS P-6b    deviation 3.997e-15 (tolerance 1.0e-12): the frame f'_p = g-tilde f_p g-tilde^dag of the project is Lord's change of representation (1.22): its Dirac matrices (3.6) are Lord's gamma (4.5) of sigma'^p = S sigma^p S^dag and equal S_LL(g) gamma^p S_LL(g)^{-1}
PASS P-7     deviation 1.554e-15 (tolerance 1.0e-12): Lord's Lorentz matrix (1.8) of L = M_LL(g-tilde) is the transpose of the project's L_LL(g) (2.10): Lambda_nu^mu = (L_LL(g))_{mu nu}
PASS P-8     deviation 1.162e-15 (tolerance 1.0e-13): the Lorentz part of the GEE derivative of the earlier hypotheses, nabla_p = d_p + Gamma^LL_p with Gamma'^LL = S_LL Gamma^LL S_LL^{-1} - (d_p S_LL) S_LL^{-1}, is covariant, and it is Lord's derivative (4.13) with Gamma_a = -Gamma^LL_a and his inhomogeneous term corrected (L4-2) (relative deviation, 12 points)

Notes:
  tolerances: exact identities of the matrices 0 or 1e-15; products of random matrices 1e-12 to 1e-14 (rounding); findings about an equation as printed pass when the printed form fails by the stated amount

VERDICT: ALL PASS (45 of 45)
\end{verbatim}

\textbf{The report of the Python program}
(\texttt{data/hck\_lg\_lord\_checks\_python.txt}):

\begin{verbatim}
PROJECT spacetimes, Revision_Lord_Gravity: lg_lord.py -- the independent check, in Python, of the Rust program lg_lord (the equations of E. A. Lord (1971) and their association with the SL(2,C) of the project)
Each line: PASS or FAIL, the identifier of the check (YL-n; the counterparts in the Rust report are named in the text),
the measured deviation and its tolerance (or the facts that decide it), and the statement.

PASS YL-1    deviation 0.000e+00 (tolerance 1.0e-13): Section 1: (1.2), (1.3), (1.5), (1.7), (1.10) to (1.15) and (1.17) of Lord hold for his matrices [L1-1 to L1-4, L1-7 to L1-10]
PASS YL-2    deviation 5.330e-15 (tolerance 1.0e-12): (1.8) to first order with L = 1 + 1/4 lambda_{mu nu} sigma^{mu nu}; for 50 random L the Lambda of (1.8) is a Lorentz matrix and satisfies (1.16) [L1-5, L1-6a, L1-6c]
PASS YL-2b   the Lambda of (1.8) is proper orthochronous [L1-6b] [50 random L]
PASS YL-3    deviation 2.665e-14 (tolerance 1.0e-12): (1.22), (1.23): epsilon S = (S^T)^{-1} epsilon; sigma-bar of S sigma S^dag is (S^dag)^{-1} sigma-bar S^{-1}; the transformed sets satisfy (1.3); for any Hermitian set (1.3) holds with N(sigma^a, sigma^b) [L1-11a to L1-13]
PASS YL-3b   a generic Hermitian set violates (1.3): Lord's 'any linearly independent set' means 'any set S sigma^mu S^dag' [L1-14] [smallest deviation from eta over 20 sets: 1.192]
PASS YL-4    deviation 7.776e-14 (tolerance 1.0e-09): (4.2) to (4.5): gamma^a(x) = S4 gamma^a S4^{-1}; (4.13), (4.31) corrected: D'_a(T psi) = T D_a psi with Gamma'_a = T Gamma_a T^{-1} + (d_a T) T^{-1} (relative, finite differences) [L4-1, L4-2]
PASS YL-4b   (4.31) as printed (T d_a T^{-1} read as a product) violates the covariance of d_a - Gamma_a [L4-3] [smallest relative violation: 0.202]
PASS YL-5    deviation 4.926e-15 (tolerance 1.0e-09): (4.32) corrected: delta a_abc = -(1/g) tr_N[gamma'_bc (d_a T) T^{-1}] [L4-4]
PASS YL-5b   (4.32) as printed (with tr_N) has the opposite sign of the inhomogeneous term [L4-5] [smallest |printed - true| / |true|: 2.00000000]
PASS YL-6    deviation 8.619e-13 (tolerance 1.0e-09): (4.19) solved: C_ba^e = 2 g a_ab^e + tr_N[(d_a gamma_b) gamma^e] (real); (4.24) C_bae = -C_eab [L4-6a, L4-6b, L4-7]
PASS YL-6b   deviation 1.610e-15 (tolerance 1.0e-09): (4.20) as printed with tr = tr_N is exactly half of that C [L4-6c]
PASS YL-6c   (4.20) as printed with the full trace is not C [L4-6d] [smallest relative deviation of the whole array: 0.5926]
PASS YL-7    deviation 3.852e-14 (tolerance 1.0e-08): (4.48): [D_mu, D_nu] psi = (d_nu Gamma_mu - d_mu Gamma_nu + [Gamma_mu, Gamma_nu]) psi (nested finite differences) [L4-9]
PASS YL-7b   deviation 7.006e-15 (tolerance 1.0e-09): the spinor image Gamma_mu = 1/4 omega_{mu ab} gamma^{ab} of a vector connection keeps gamma_b constant, and its curvature is 1/4 R_{mu nu alpha beta} gamma^{alpha beta} [L4-10a, L4-10b]
PASS YL-7c   the printed coefficient 1/2 of (4.47), (4.48) gives twice the spinor curvature [L4-10c] [smallest relative deviation: 1.00000000]
PASS YL-8    deviation 4.761e-13 (tolerance 1.0e-11): the dictionary: Lord's sigma, sigma-bar, gamma, C equal the project's; his gamma^5 = -gamma^ch_LL; his diag(S, (S^dag)^{-1}) with S = M_LL(g-tilde) is S_LL(g); M_LL(b-bar) = -epsilon M_LL(b)^T epsilon; the frame g-tilde f g-tilde^dag has Lord's Dirac matrices; his Lambda of M_LL(g-tilde) is L_LL(g)^T [P-1 to P-7]

Notes:
  derivatives by central finite differences of eighth order (step 0.01), nested for second derivatives; tolerances 1e-12 for exact matrix identities, 1e-8 to 1e-9 where finite differences enter

VERDICT: ALL PASS (16 of 16)
\end{verbatim}

\textbf{The report of the Wolfram script}
(\texttt{data/hck\_lg\_lord\_checks\_wolfram.txt}):

\begin{verbatim}
lg_lord.wls: the exact identities and factors of Lord (1971), Sections 1 and 4 (Revision_Lord_Gravity, PROJECT spacetimes)
Each line: PASS or FAIL, the identifier of the check in this script (WL and a number), the statement and, in brackets, the
facts that decide it.  Exact arithmetic: symbolic matrices, every identity simplified to zero.

PASS WL-1     (Rust L1-1 to L1-4, L1-7 to L1-10; Python YL-1) Lord's identities (1.2), (1.3), (1.5) (256 index combinations), (1.7), (1.10), (1.11), (1.12), (1.13) (both forms), (1.14), (1.15) and (1.17) hold exactly for his matrices [(1.2) True, (1.3) True, (1.5) True, (1.7) True, (1.10)-(1.14) True, (1.15) True, (1.17) True]
PASS WL-2     (Rust L1-11a, L1-11b; Python YL-3) for S of determinant 1 (symbolic entries): epsilon S = (S^T)^{-1} epsilon; and for any matrix sigma and a second S' of determinant 1, bar(S sigma S') = S'^{-1} bar(sigma) S^{-1}; with S' = S^dag this is (1.23) [epsilon True, (1.23) True]
PASS WL-3     (Rust L1-13, L1-14; Python YL-3, YL-3b) for any two 2x2 matrices A, B (symbolic entries): 1/2 (A bar(B) + B bar(A)) = N(A, B) 1 with N(A, B) = (det(A + B) - det A - det B)/2, and N(sigma^mu, sigma^nu) = eta^{mu nu}; so (1.3) holds for a set of four matrices exactly when their products N are eta^{ab} [identity True, N of the standard set True]
PASS WL-4     (Rust L1-5; Python YL-2) (1.8) to first order: with Lambda_mu^nu = delta_mu^nu + eta^{nu rho} lambda_{rho mu} and L = 1 + 1/4 lambda_{mu nu} sigma^{mu nu} (six real symbols), eta^{mu rho} lambda_{rho nu} sigma^nu + X sigma^mu + sigma^mu X^dag = 0 for every mu [True]
PASS WL-5     (Rust L4-2, L4-3; Python YL-4, YL-4b) for S(t) = ((1 + t^3, t), (t^2, 1)) (determinant 1), a constant Gamma = diag(Y, -Y^dag) of the Lorentz algebra and a spinor of four symbolic functions: D = d - Gamma is covariant with Gamma' = S Gamma S^{-1} + (dS) S^{-1}, and not with the printed inhomogeneous term S (d S^{-1}) [corrected law covariant True, printed law not covariant True]
PASS WL-6     (Rust L4-6b, L4-6c; Python YL-6, YL-6b) in Lord's special gauge (constant gamma, (5.2)): for Gamma = 1/2 g a_bc gamma^{bc} with six symbolic a_bc, [Gamma, gamma_b] = -2 g a_b^e gamma_e, so (4.19) gives C_b^e = 2 g a_b^e: twice the printed (4.20) and (5.2), C = g a [C = 2 g a: True, in the span of gamma_e: True]
PASS WL-7     (Rust L4-10a, L4-10b, L4-10c; Python YL-7b, YL-7c) lambda(W) = 1/4 W_{ab} gamma^{ab} (W antisymmetric, symbolic) satisfies [lambda(W), gamma_b] = -W_b^e gamma_e and [lambda(W1), lambda(W2)] = lambda([W1, W2]): a homomorphism of Lie algebras, so the curvature of the spinor connection lambda(W_mu) is lambda(R_{mu nu}) = 1/4 R_{mu nu ab} gamma^{ab} exactly, and Lord's printed 1/2 R gamma is twice it [action on gamma True, homomorphism True]

Notes:
  Wolfram Language 15.0.1; run time 0.0 s

VERDICT: ALL PASS (7 of 7)
\end{verbatim}

\hypertarget{the-symbols-of-this-manuscript}{%
\section{10. The symbols of this
manuscript}\label{the-symbols-of-this-manuscript}}

\begin{longtable}[]{@{}
  >{\raggedright\arraybackslash}p{(\columnwidth - 2\tabcolsep) * \real{0.5000}}
  >{\raggedright\arraybackslash}p{(\columnwidth - 2\tabcolsep) * \real{0.5000}}@{}}
\toprule\noalign{}
\begin{minipage}[b]{\linewidth}\raggedright
symbol
\end{minipage} & \begin{minipage}[b]{\linewidth}\raggedright
meaning
\end{minipage} \\
\midrule\noalign{}
\endhead
\bottomrule\noalign{}
\endlastfoot
\(\eta^{\rm LL}_{\mu\nu}\), \(\eta^{{\rm LL}\,\mu\nu}\) & the Minkowski
metric \(\mathrm{diag}(1,-1,-1,-1)\) and its inverse \\
\(\sigma^{\rm LL}_\mu\) & the unit matrix and the Pauli matrices (Lord's
\(\sigma^\mu\)) \\
\(\bar X\) & the bar of a \(2\times2\) matrix,
\(-\mathsf E_{\rm LL}X^T\mathsf E_{\rm LL}\), its adjugate (1.1) \\
\(\mathsf E_{\rm LL}\) & the matrix
\(\begin{pmatrix}0&1\\-1&0\end{pmatrix}\) (Lord's \(\epsilon\)) \\
\(\mathsf 1_2\), \(\mathsf 1_4\) & the unit matrices \\
\(\gamma^\mu_{\rm LL}\), \(\gamma^{\mu\nu}_{\rm LL}\) & the Dirac
matrices in the Weyl form (the chiral basis) and their antisymmetrized
products \\
\(\gamma^{\rm ch}_{\rm LL}\) & the chirality matrix
\(i\gamma^0_{\rm LL}\gamma^1_{\rm LL}\gamma^2_{\rm LL}\gamma^3_{\rm LL}\)
(minus Lord's \(\gamma^5\)) \\
\(C_{\rm LL}\) & the charge-conjugation matrix
\(\mathrm{diag}(\mathsf E_{\rm LL},-\mathsf E_{\rm LL})\) (Lord's
\(C\)) \\
\(\sigma^{\mu\nu}_{\rm LG}\) &
\(\tfrac12(\sigma^{\rm LL}_\mu\bar\sigma^{\rm LL}_\nu-\sigma^{\rm LL}_\nu\bar\sigma^{\rm LL}_\mu)\)
(Lord's \(\sigma^{\mu\nu}\)) \\
\(\epsilon^{\rm LG}_{\mu\nu\rho\sigma}\) & the Levi-Civita symbol,
\(\epsilon^{\rm LG}_{0123}=+1\) \\
\(f_{\bar\mu\bar\nu}\) & the dual
\(\tfrac12i\,\epsilon^{\rm LG}_{\mu\nu\rho\sigma}f^{\rho\sigma}\) of an
antisymmetric tensor \\
\(\mathrm{tr}\), \(\mathrm{tr}_{\rm N}\) & the trace, and the trace of a
\(4\times4\) matrix normalized to
\(\mathrm{tr}_{\rm N}\mathsf 1_4=1\) \\
\(N_{\rm LG}(X,Y)\) & the polarized determinant
\(\tfrac12(\det(X+Y)-\det X-\det Y)\), the Minkowski product of two
Hermitian matrices (2.7) \\
\(\Lambda_{\rm Lord}\), \(\lambda^{\rm Lord}_{\mu\nu}\) & a Lorentz
matrix of Lord's (2.3), and the parameters of an infinitesimal one \\
\(L_{\rm Lord}\), \(\mathsf L_{\rm Lord}\) & its spinor matrices of
\(SL(2,\mathbb C)\) (\(2\times2\)) and on Dirac spinors
(\(4\times4\)) \\
\(\psi^{\rm Lord}_{A_1\dots A_{2s}}\) & a Bargmann-Wigner field of spin
\(s\); capital Latin letters are Dirac spinor indices \\
\(m\) & the rest mass of a Bargmann-Wigner field \\
\(S_{\rm Lord}\), \(\mathsf S_{\rm Lord}\) & a change of representation
(2.6): a unimodular \(2\times2\) matrix, and its \(4\times4\) form
(5.1) \\
\(\mathsf T\) & a matrix field by which a spinor index is transformed,
in (5.5) and (5.6) \\
\(\sigma^a_{\rm Lord}(x)\), \(\gamma^a_{\rm Lord}(x)\) & the Pauli-type
matrices and the Dirac matrices as fields (5.2) \\
\(d^{\rm Lord}_a\), \(X^{\rm Lord}_{ab}{}^c\) & Lord's derivative
operators and the coefficients of their commutators \\
\(D^{\rm Lord}_a\) & Lord's generalized derivative (5.3) \\
\(\Gamma^{\rm Lord}_a\), \(a^{\rm Lord}_{abc}\), \(g_{\rm Lord}\) & the
auxiliary fields of the spinors, their coefficients (5.4), and the
universal coupling \\
\(C^{\rm Lord}_{ba}{}^e\) & the connection of the vectors (5.7),
(5.8) \\
\(h^{\rm Lord}_a{}^\mu\), \(g^{\rm Lord}_{\mu\nu}\) & the reference
vectors (tetrad) and the metric \\
\(\Gamma^{\rm Lord}_{\mu\nu}{}^\rho\), \(T^{\rm Lord}_{\mu\nu}{}^\rho\),
\(R^{\rm Lord}_{\mu\nu\rho}{}^\sigma\) & the affine connection, its
torsion and its curvature (5.9) \\
\(W_\mu\), \(W_{\mu\,ab}\) & a connection of the vectors as a matrix
\((W_\mu)_b{}^e\) and with both indices lowered (5.11) \\
\(R_{\mu\nu\alpha\beta}\) & the curvature of a connection of the
vectors, with its last two indices lowered (5.11) \\
\(A^{\rm Lord}_{\mu\nu\rho\sigma}\),
\(A^{\rm Lord}_{\rho\sigma\lambda}\), \(A^{\rm Lord}_{\mu\nu}\) & the
real spin-2 tensor, its potentials and its secondary potentials \\
\(B^{\rm Lord}_{\mu\nu\rho\sigma}\) & the tensor of Lord's Section 3
that has the symmetries of the Riemann tensor \\
\(j^{\rm Lord}_{\nu\rho\sigma}\), \(j^{\rm Lord}_{\mu\nu}\),
\(j^{\rm Lord}\), \(f_{\rm Lord}\) & the current of the spin-2 field,
the symmetric tensor from which it derives, its trace, and the coupling
of Lord's Section 3 \\
\(f_{\mu\nu}\), \(A_\mu\), \(e\) & the electromagnetic field strength,
its potential and the electric charge (Lord's Sections 1 and 2) \\
\end{longtable}

\end{document}
